\section{Типология полярных сияний и определяющие факторы}

\subsection{Диффузное и дискретное сияние}

Выделяют два главных типа полярных сияний по форме и физическим механизмам (см. табл.~\ref{tab:types}).

\begin{table}[H]
    \centering
    \caption{Сравнение диффузного и дискретного сияний}
    \label{tab:types}
    \begin{tabular}{|p{4cm}|p{5cm}|p{5cm}|}
        \hline
        \textbf{Характеристика} & \textbf{Диффузное сияние} & \textbf{Дискретное сияние} \\
        \hline
        Форма & Равномерное свечение & Чёткие структуры (дуги, лучи) \\
        \hline
        Размер & Сотни километров & 1--10 км \\
        \hline
        Длительность & Несколько минут & Секунды \\
        \hline
        Механизм ускорения & Альфвеновский & Квазистатический \\
        \hline
    \end{tabular}
\end{table}

\subsection{Зависимость цвета от высоты и энергии частиц}

Цвет свечения зависит от состава атмосферы и энергии электронов:
\begin{itemize}
    \item Электроны $\sim$1 кэВ вызывают красное свечение ($\lambda = 630$ нм) на высоте $\sim200$ км.
    \item Электроны 5--10 кэВ создают зелёное свечение ($\lambda = 557,7$ нм) на высоте 100--120 км.
    \item Электроны $>$30 кэВ проникают ниже 100 км и вызывают синие и фиолетовые оттенки ($\lambda = 391,4$ нм и $427,8$ нм).
\end{itemize}

\subsection{Внешние факторы}

Интенсивность и частота полярных сияний зависят от:
\begin{enumerate}
    \item \textbf{11-летнего солнечного цикла:} в максимуме активности сияния ярче и чаще.
    \item \textbf{Геомагнитной активности (индекс Kp):} при Kp $\ge$ 5 наблюдаются бури, сияния видны в средних широтах.
    \item \textbf{Сезона:} максимум активности приходится на март и сентябрь (равноденствия).
\end{enumerate}